\section{Verification}\label{sup:verification}

The code of the paper is a Python package with a test suite in four families:
\emph{construction} (embeddings, multiplication tables, inverses, the hexagonal extra
modulus and the changes of basis), \emph{sources} (comparisons with published results
and controls of the transcribed data), \emph{averages and application} (orientation
averages, concentration tensor, averaged Eshelby tensor and self-consistent scheme) and
the \emph{manuscript chain} (every computed number, table and figure of the paper and of
this supplement is regenerated from the code and compared with the files used, and the
formulas printed in sections~2--6 and appendix~A of the paper are restated in the tests as
printed and checked).

Each comparison has one of three levels. \emph{Verification} (V) is against exact
mathematics, computed with exact rational and algebraic arithmetic, or against an
independent route of our own: brute force over all products, inversion or spectrum of
the $18\times18$ matrix of the tensor on $\Vsg$, random rotations. \emph{Cross-comparison}
(CC) is against a published result, reconstructed from the paper that gives it; where a
printed formula does not hold, the test checks both the printed and the corrected form.
\emph{Consistency} (C) covers limit cases (isotropic and separable tensors, isotropic
grains), invariances and inequalities checked on random samples. No comparison is with
measurements: the inputs are published atomistic constants.

Table~\ref{tab:s-audit} maps every claim of the paper, and the illustrative application,
to the key tests that check it. These forty-two tests are marked \texttt{core} and run with
\texttt{pytest -m core} in about a minute and a half, including the exact and the slow
ones; a test of the suite fails if the table names a test that is not marked, or if a
marked test is missing from the table. The default \texttt{pytest} runs the fast
numerical tests, and \texttt{pytest -{}-all} the whole suite in a few minutes.

\newcommand{\AuditMapCaption}{Audit map of the verification. Tests are named as
\texttt{file:function}, for the function \texttt{test\_function} of
\texttt{tests/test\_file.py} in the code repository. Levels: V, verification; CC,
cross-comparison with a published result; C, consistency.\label{tab:s-audit}}
\begingroup\footnotesize\setlength{\LTcapwidth}{\linewidth}
\input{audit_map}
\endgroup
